%Cấu hình tên các tham số
%
Tên tham số 1: Lớp
Tên tham số 2: Môn
Tên tham số 3: Chương
Tên tham số 4: Mức độ
Tên tham số 5: Bài
Tên tham số 6: Dạng
%
%Cấu hình chi tiết ID
%
%%Cấu hình mức độ dùng chung.
[Y] Yếu
[B] Trung bình
[K] Khá
[G] Giỏi
[T] Thực tế
[N] Nhận biết
[H] Thông hiểu
[V] Vận dụng
[C] Vận dụng cao

-[8] Lớp 8

----[K] Khoa học tự nhiên

-------[1] Phản ứng hóa học

----------[1] Biến đổi vật lí và biến đổi hóa học
-------------[1] Tìm hiểu và phân biệt sự biến đổi vật lí, sự biến đổi hóa học
-------------[2] Bài tập nhận biết hiện tượng hóa học và hiện tượng vật lí trong tự nhiên
-------------[3] Giải thích các hiện tượng thực tiễn: nước đá tan, đốt cháy than, rỉ sét...

----------[2] Phản ứng hóa học và năng lượng
-------------[1] Lý thuyết về diễn biến và dấu hiệu của phản ứng hóa học
-------------[2] Phân biệt phản ứng tỏa nhiệt và thu nhiệt
-------------[3] Ứng dụng của phản ứng tỏa nhiệt trong đời sống (đốt gas, than) và sản xuất

----------[3] Định luật bảo toàn khối lượng và phương trình hóa học
-------------[1] Phát biểu định luật và viết phương trình bảo toàn khối lượng
-------------[2] Lập phương trình hóa học từ sơ đồ phản ứng
-------------[3] Bài toán tính khối lượng chất tham gia và sản phẩm dựa vào định luật bảo toàn khối lượng

----------[4] Mol và tỉ khối của chất khí
-------------[1] Khái niệm mol, khối lượng mol và thể tích mol
-------------[2] Bài tập chuyển đổi giữa số mol, khối lượng và thể tích
-------------[3] Bài tập tính tỉ khối của chất khí
-------------[4] Ứng dụng thực tiễn: giải thích tại sao bóng bay bơm khí hydro/heli bay lên được

----------[5] Tính theo phương trình hóa học
-------------[1] Tính khối lượng, số mol, thể tích của chất phản ứng và sản phẩm
-------------[2] Bài toán tính hiệu suất phản ứng
-------------[3] Bài toán lượng dư trong phản ứng hóa học

----------[6] Nồng độ dung dịch
-------------[1] Khái niệm độ tan và các yếu tố ảnh hưởng đến độ tan
-------------[2] Bài toán tính nồng độ phần trăm (C%) và nồng độ mol (CM)
-------------[3] Bài toán pha chế dung dịch theo nồng độ cho trước
-------------[4] Bài tập thực tiễn: tính nồng độ muối/đường trong ẩm thực, y tế (nước muối sinh lí)

----------[7] Tốc độ phản ứng và chất xúc tác
-------------[1] Khái niệm và các yếu tố ảnh hưởng đến tốc độ phản ứng
-------------[2] Vai trò của chất xúc tác và chất ức chế
-------------[3] Bài tập vận dụng giải thích các biện pháp bảo quản thực phẩm (bơm khí, hút chân không, bảo quản lạnh)

-------[2] Một số hợp chất thông dụng

----------[1] Acid
-------------[1] Khái niệm và tính chất hóa học của acid
-------------[2] Bài toán định lượng liên quan đến acid
-------------[3] Ứng dụng của một số acid thông dụng (HCl, H2SO4, CH3COOH) trong thực tiễn

----------[2] Base
-------------[1] Khái niệm, phân loại và tính chất hóa học của base
-------------[2] Bài toán định lượng liên quan đến base
-------------[3] Ứng dụng của base trong công nghiệp và đời sống (Ca(OH)2 khử chua đất...)

----------[3] Thang pH
-------------[1] Khái niệm, ý nghĩa của thang pH
-------------[2] Xác định môi trường dung dịch dựa vào pH
-------------[3] Bài tập thực tiễn: vai trò của pH đối với đất nông nghiệp, nước sinh hoạt, dạ dày con người

----------[4] Oxide
-------------[1] Khái niệm, phân loại oxide
-------------[2] Tính chất hóa học của các loại oxide
-------------[3] Bài toán nhận biết và tách oxide

----------[5] Muối
-------------[1] Khái niệm, phân loại và tính tan của muối
-------------[2] Tính chất hóa học và phương pháp điều chế muối
-------------[3] Bài toán về muối (tác dụng với acid, base, muối khác)
-------------[4] Mối liên hệ giữa các hợp chất vô cơ

----------[6] Phân bón hóa học
-------------[1] Thành phần, vai trò của các loại phân bón hóa học (đạm, lân, kali)
-------------[2] Cách sử dụng phân bón hóa học an toàn, hiệu quả
-------------[3] Bài tập tính hàm lượng dinh dưỡng trong phân bón
-------------[4] Tác động của phân bón đến môi trường đất và nước sinh hoạt

-------[3] Khối lượng riêng và áp suất

----------[1] Khối lượng riêng
-------------[1] Định nghĩa, công thức tính khối lượng riêng
-------------[2] Bài tập tính toán khối lượng riêng, khối lượng, thể tích
-------------[3] Thực hành xác định khối lượng riêng của chất rắn/lỏng bằng dụng cụ đo
-------------[4] Giải thích vì sao dầu nổi trên nước, phân loại vật liệu dựa vào khối lượng riêng

----------[2] Áp suất chất rắn
-------------[1] Áp lực và công thức tính áp suất
-------------[2] Bài tập tính toán áp suất trên một mặt phẳng
-------------[3] Giải thích thực tiễn: tại sao dao cần mài sắc, xe tăng dùng bánh xích thay vì bánh xe

----------[3] Áp suất chất lỏng và chất khí
-------------[1] Đặc điểm áp suất chất lỏng, bình thông nhau
-------------[2] Sự truyền áp suất chất lỏng (nguyên lí Pascal) và máy nén thủy lực
-------------[3] Áp suất khí quyển và sự phụ thuộc vào độ cao
-------------[4] Bài tập thực tiễn: nguyên lí bình xịt, máy đo huyết áp, giác hút

----------[4] Lực đẩy Archimedes
-------------[1] Lực đẩy của chất lỏng lên vật đặt trong nó (Lực đẩy Archimedes)
-------------[2] Điều kiện định tính để một vật nổi, chìm hay lơ lửng
-------------[3] Bài tập tính toán lực đẩy Archimedes
-------------[4] Bài tập thực tiễn: giải thích sự nổi của tàu thuyền, khinh khí cầu, phao bơi

-------[4] Tác dụng làm quay của lực

----------[1] Moment lực
-------------[1] Tác dụng làm quay của lực, khái niệm moment lực
-------------[2] Các yếu tố ảnh hưởng đến moment lực
-------------[3] Giải thích thực tiễn: tại sao tay nắm cửa đặt ở mép xa bản lề, cách vặn cờ lê dễ dàng

----------[2] Đòn bẩy
-------------[1] Khái niệm và các loại đòn bẩy
-------------[2] Tác dụng làm thay đổi hướng và độ lớn của lực
-------------[3] Ứng dụng đòn bẩy trong cơ thể người (khớp xương) và công cụ lao động (kéo, kìm, xe cút kít)

-------[5] Điện

----------[1] Sự nhiễm điện do cọ xát
-------------[1] Hiện tượng nhiễm điện do cọ xát
-------------[2] Giải thích nguyên nhân các vật bị nhiễm điện
-------------[3] Liên hệ thực tiễn: hiện tượng sấm sét, tóc dựng đứng mùa hanh khô, bụi bám cánh quạt

----------[2] Dòng điện và mạch điện
-------------[1] Khái niệm dòng điện, nguồn điện, vật dẫn điện và vật cách điện
-------------[2] Sơ đồ mạch điện đơn giản, kí hiệu các thiết bị điện
-------------[3] Thực hành lắp mạch điện cơ bản (nguồn, công tắc, bóng đèn)
-------------[4] Bài tập nhận biết lỗi sai trong mạch điện

----------[3] Tác dụng của dòng điện
-------------[1] Tác dụng nhiệt và tác dụng phát sáng
-------------[2] Tác dụng từ, hóa học và sinh lí
-------------[3] Giải thích ứng dụng các tác dụng của dòng điện (mạ điện, chuông điện, bóng đèn LED)
-------------[4] An toàn điện: tác động của dòng điện lên cơ thể, quy tắc phòng tránh

-------[6] Nhiệt

----------[1] Năng lượng nhiệt và nội năng
-------------[1] Khái niệm năng lượng nhiệt, nội năng
-------------[2] Mối quan hệ giữa nhiệt độ và nội năng của vật
-------------[3] Giải thích sự truyền nhiệt giữa các vật khi tiếp xúc

----------[2] Sự truyền nhiệt
-------------[1] Dẫn nhiệt: đặc điểm, chất dẫn nhiệt tốt và kém
-------------[2] Đối lưu: quá trình truyền nhiệt trong chất lỏng và chất khí
-------------[3] Bức xạ nhiệt: truyền nhiệt qua chân không
-------------[4] Giải thích thực tiễn: áo ấm mùa đông, phích nước, nồi xoong dùng kim loại cán nhựa

----------[3] Sự nở vì nhiệt
-------------[1] Đặc điểm sự nở vì nhiệt của chất rắn, lỏng, khí
-------------[2] Ứng dụng sự nở vì nhiệt trong thực tiễn (nhiệt kế, rơ le nhiệt)
-------------[3] Tác hại và biện pháp khắc phục sự nở vì nhiệt (khe hở đường ray, cầu thép)

-------[7] Sinh học cơ thể người

----------[1] Khái quát về cơ thể người
-------------[1] Cấu tạo cơ bản của cơ thể người, các hệ cơ quan
-------------[2] Mối quan hệ phối hợp hoạt động giữa các hệ cơ quan

----------[2] Hệ vận động
-------------[1] Cấu tạo và chức năng của bộ xương, hệ cơ
-------------[2] Sự phù hợp giữa cấu tạo và chức năng của xương, khớp, cơ vân
-------------[3] Bài tập tình huống: sơ cứu gãy xương, băng bó
-------------[4] Các bệnh liên quan đến hệ vận động (loãng xương, cong vẹo cột sống) và cách phòng tránh

----------[3] Dinh dưỡng và tiêu hóa
-------------[1] Chức năng của các cơ quan trong hệ tiêu hóa
-------------[2] Quá trình tiêu hóa và hấp thụ chất dinh dưỡng
-------------[3] Chế độ dinh dưỡng hợp lí, an toàn vệ sinh thực phẩm
-------------[4] Các bệnh tiêu hóa (đau dạ dày, ngộ độc thực phẩm) và cách bảo vệ

----------[4] Máu và hệ tuần hoàn
-------------[1] Thành phần và chức năng của máu, miễn dịch và vaccine
-------------[2] Các nhóm máu ở người và nguyên tắc truyền máu
-------------[3] Cấu tạo tim, mạch máu và chu kì tuần hoàn
-------------[4] Thực hành sơ cứu cầm máu, đo huyết áp
-------------[5] Bệnh tim mạch (huyết áp cao, nhồi máu cơ tim, đột quỵ) và phòng tránh

----------[5] Hệ hô hấp
-------------[1] Cấu tạo và chức năng hệ hô hấp
-------------[2] Quá trình trao đổi khí ở phổi và tế bào
-------------[3] Tác hại của khói thuốc lá và ô nhiễm không khí đối với hô hấp
-------------[4] Thực hành hô hấp nhân tạo cho người đuối nước/ngạt khí

----------[6] Hệ bài tiết
-------------[1] Cấu tạo hệ bài tiết nước tiểu
-------------[2] Quá trình hình thành nước tiểu
-------------[3] Một số bệnh hệ bài tiết (sỏi thận, suy thận) và thói quen sinh hoạt bảo vệ

----------[7] Điều hòa môi trường trong cơ thể
-------------[1] Khái niệm nội môi, vai trò của điều hòa môi trường trong
-------------[2] Cơ chế điều hòa thân nhiệt
-------------[3] Liên hệ thực tiễn: cảm nắng, cảm lạnh, biện pháp phòng chống
-------------[4] Sự điều hòa lượng đường trong máu

----------[8] Hệ thần kinh và giác quan
-------------[1] Cấu tạo và chức năng hệ thần kinh (não, tủy sống, dây thần kinh)
-------------[2] Phản xạ, cung phản xạ
-------------[3] Cấu tạo cơ quan thị giác (mắt) và thính giác (tai)
-------------[4] Các bệnh về mắt (cận thị, viễn thị) và tai, cách bảo vệ
-------------[5] Tác hại của chất gây nghiện, ma túy đối với hệ thần kinh

----------[9] Hệ nội tiết
-------------[1] Vai trò của các tuyến nội tiết (tuyến yên, tuyến giáp, tuyến tụy, tuyến sinh dục)
-------------[2] Một số bệnh nội tiết phổ biến (tiểu đường, bướu cổ)
-------------[3] Sự thay đổi cơ thể tuổi dậy thì và vệ sinh sinh dục

-------[8] Sinh thái

----------[1] Môi trường và các nhân tố sinh thái
-------------[1] Khái niệm môi trường sống, các loại môi trường
-------------[2] Nhân tố sinh thái vô sinh và hữu sinh
-------------[3] Giới hạn sinh thái và sự thích nghi của sinh vật
-------------[4] Phân tích tác động của ánh sáng, nhiệt độ, độ ẩm lên sinh vật tại địa phương

----------[2] Quần thể sinh vật
-------------[1] Khái niệm quần thể sinh vật, dấu hiệu nhận biết
-------------[2] Các đặc trưng cơ bản (mật độ, tỉ lệ giới tính, cấu trúc tuổi)
-------------[3] Liên hệ thực tiễn: ý nghĩa của việc duy trì mật độ hợp lí trong nông nghiệp

----------[3] Quần xã sinh vật
-------------[1] Khái niệm quần xã, các đặc trưng cơ bản
-------------[2] Mối quan hệ giữa các loài trong quần xã (cạnh tranh, cộng sinh, kí sinh...)
-------------[3] Bài tập nhận dạng mối quan hệ sinh thái trong tự nhiên

----------[4] Hệ sinh thái
-------------[1] Khái niệm và thành phần cấu trúc của hệ sinh thái
-------------[2] Chuỗi thức ăn, lưới thức ăn và tháp sinh thái
-------------[3] Phân tích chuỗi thức ăn trong một hệ sinh thái cụ thể (ao, rừng, đồng cỏ)
-------------[4] Giải thích hậu quả sinh thái khi một loài tiêu diệt loài khác (mất cân bằng)

----------[5] Cân bằng tự nhiên và bảo vệ môi trường
-------------[1] Khái niệm cân bằng tự nhiên và các nguyên nhân gây mất cân bằng
-------------[2] Tác động của con người đối với môi trường (ô nhiễm, khai thác quá mức)
-------------[3] Biện pháp bảo vệ, duy trì cân bằng tự nhiên và bảo vệ động vật hoang dã
-------------[4] Đề xuất dự án nhỏ về bảo vệ môi trường tại khu dân cư/trường học